\documentclass{article}
\usepackage[utf8]{inputenc}
\usepackage[T2A]{fontenc}
\usepackage[russian]{babel}
\usepackage{amsmath}
\usepackage{amssymb}

\title{Лекция 3 \\ Синдромное декодирование и границы параметров кодов}
\author{}
\date{}

\begin{document}

\maketitle

\section{Синдромное декодирование линейных блочных кодов}

Синдромное декодирование — это эффективный способ декодирования линейных блочных кодов, использующий проверочную матрицу. 

Постановка задачи: пусть передаётся кодовое слово \( \mathbf{c} \in C \) линейного \((n,k)\)-кода, а принимается вектор \( \mathbf{r} = \mathbf{c} + \mathbf{e} \), где \( \mathbf{e} \) — вектор ошибки. Декодеру неизвестны ни \( \mathbf{c} \), ни \( \mathbf{e} \); известно только \( \mathbf{r} \). Требуется восстановить наиболее вероятное переданное слово.

Скорость кода \( R = k/n \) показывает долю информационных символов в кодовом слове.

Синдром вычисляется как \( \mathbf{s} = \mathbf{r} \mathbf{H}^T = (\mathbf{c} + \mathbf{e}) \mathbf{H}^T = \mathbf{e} \mathbf{H}^T \), поскольку \( \mathbf{c} \mathbf{H}^T = 0 \). Синдром зависит только от ошибки.

Алгоритм декодирования:
\begin{enumerate}
    \item Вычислить синдром \( \mathbf{s} \).
    \item Найти по синдрому наиболее вероятный вектор ошибки \( \mathbf{e} \) (лидер смежного класса).
    \item Восстановить кодовое слово \( \mathbf{c} = \mathbf{r} - \mathbf{e} \).
\end{enumerate}

В двоичном случае существует \( 2^{n-k} \) различных синдромов, что соответствует числу смежных классов. Таблица синдромов сопоставляет каждому синдрому лидера смежного класса — вектор ошибки наименьшего веса в этом классе.

\textbf{Пример:} код Хэмминга \((7,4)\) имеет проверочную матрицу размера \(3 \times 7\), \( 2^3 = 8 \) синдромов, каждый указывает на однократную ошибку. Табличное декодирование не требует перебора всех кодовых слов.

\section{Радиус покрытия кода}

Радиус покрытия \( \rho(C) \) кода \( C \) в пространстве Хэмминга \( \mathbb{F}_q^n \) — это максимальное расстояние от произвольного вектора пространства до ближайшего кодового слова:
\[
\rho(C) = \max_{\mathbf{u} \in \mathbb{F}_q^n} \min_{\mathbf{v} \in C} d(\mathbf{u}, \mathbf{v}),
\]
где \( d(\cdot,\cdot) \) — расстояние Хэмминга. Это наименьший радиус \( r \), при котором шары радиуса \( r \) с центрами в кодовых словах покрывают всё пространство.

Радиус покрытия показывает, насколько далеко принятое слово может оказаться от множества допустимых слов. Для совершенных кодов, например кодов Хэмминга \( (2^r-1, 2^r-r-1) \), радиус покрытия равен 1: шары единичного радиуса не пересекаются и полностью покрывают пространство. В общем случае радиус покрытия не меньше радиуса упаковки, и равенство достигается только для совершенных кодов.

\section{Граница Хэмминга (граница сферической упаковки)}

Граница Хэмминга даёт верхнюю оценку мощности \( q \)-ичного кода длины \( n \) с минимальным расстоянием \( d \). Пусть код исправляет \( t = \lfloor (d-1)/2 \rfloor \) ошибок. Тогда шары радиуса \( t \) с центрами в кодовых словах не пересекаются. Объём такого шара:
\[
V_q(n, t) = \sum_{i=0}^{t} \binom{n}{i} (q-1)^i.
\]
Суммарный объём \( M \) непересекающихся шаров не может превышать объём всего пространства \( q^n \). Отсюда:
\[
M \le \frac{q^n}{\sum_{i=0}^{t} \binom{n}{i} (q-1)^i}.
\]

\textbf{Теорема:} для любого \( q \)-ичного кода длины \( n \) с расстоянием \( d > 2t+1 \) число кодовых слов \( M \) удовлетворяет этому неравенству. Доказательство основано на непересекаемости шаров. Коды, достигающие границы, называются совершенными (коды Хэмминга, коды Голея, тривиальные коды).

\section{Граница Варшамова — Гильберта}

Граница Варшамова — Гильберта — это нижняя оценка существования линейных кодов с заданными параметрами.

\textbf{Теорема:} если для заданных \( n, k, d \) и размера алфавита \( q \) выполняется неравенство
\[
q^{n-k} > \sum_{i=0}^{d-2} \binom{n-1}{i} (q-1)^i,
\]
то существует линейный \( q \)-ичный \( (n, k) \)-код с минимальным расстоянием не менее \( d \).

\textbf{Идея доказательства:} столбцы проверочной матрицы \( H \) размера \( (n-k) \times n \) выбираются последовательно так, чтобы любые \( d-1 \) столбцов были линейно независимы. На каждом шаге новый столбец не должен быть линейной комбинацией \( d-2 \) или меньшего числа уже выбранных столбцов. Количество «запрещённых» столбцов на \( j \)-м шаге не превышает \( \sum_{i=0}^{d-2} \binom{j-1}{i} (q-1)^i \). Если общее число возможных столбцов \( q^{n-k} \) больше этого количества при \( j = n \), построение возможно.

В асимптотике граница принимает вид:
\[
R \ge 1 - H_q(\delta),
\]
где \( R = k/n \) — скорость кода, \( \delta = d/n \) — относительное расстояние, \( H_q \) — \( q \)-ичная энтропийная функция. Граница гарантирует существование кода, но не даёт явного эффективного построения.

\end{document}